\documentclass[10pt,a4paper]{amsart}
\usepackage[margin=19mm]{geometry}
\usepackage{amsmath,amssymb,mathtools}
\usepackage[scr=rsfs,cal=euler]{mathalfa}
\usepackage{xfrac}
\usepackage[unicode,hidelinks,bookmarks=false]{hyperref}
\newcommand{\ie}{i.e.\ }
\newcommand{\cf}{cf.\ }
\newcommand{\resp}{respectively\ }
\newcommand{\Subsection}{\subsection}
\providecommand{\nobreakdash}{\nobreak\hbox{-}}
\setlength{\emergencystretch}{2em}
\multlinegap0pt
\usepackage{amssymb}
\usepackage{mathtools}
\usepackage[shortlabels]{enumitem}
\newtheorem{theorem}{Theorem}
\newtheorem{corollary}{Corollary}
\newcommand{\Bt}{B^2}
\newcommand{\Fp}{\mathbb{F}_p}
\newcommand{\Ht}{H^2}
\newcommand{\Zt}{Z^2}
\newcommand{\kk}{\Bbbk}
\title{Cohomology and extensions of Novikov algebras of truncated polynomials}
\author{Hassan Alhussein}
\date{Source: arXiv:2608.25372v1 (CC BY 4.0)}
\begin{document}
\maketitle

\noindent\emph{Source and licence.} Cohomology and extensions of Novikov algebras of truncated polynomials. Version arXiv:2608.25372v1 (CC BY 4.0), distributed by arXiv under the Creative Commons Attribution 4.0 International licence, http://creativecommons.org/licenses/by/4.0/, as recorded in the arXiv licence metadata for this exact version and shown on the abstract page https://arxiv.org/abs/2608.25372v1. Full licence notice: \texttt{licenses/arxiv-2608.25372-CC-BY-4.0.txt}.\par\smallskip

\noindent\textit{Editorial scope.} Counted unit: the article's theorem thm:p2 (source0.tex lines 965--985). The article's complete original proof of it is delivered with it (source0.tex lines 987--1042, one \texttt{proof} environment of 56 lines). No mathematics is written, repaired or re-derived: the statement and the proof are the article's own lines, in the article's own notation, and no retained line is substituted or re-flowed.  Retained uncounted as the source-local context the proof needs: lines 503--510; lines 647--649; lines 725--729.  Nothing was omitted from any retained span, so the package carries no omission record.  No retained line contains a citation, so the wrapper carries no bibliography and no bibliography entry.  No retained line contains a figure environment, an image-inclusion command or drawing code, so no image is delivered and the package declares no asset dependency.  Editorial formatting only, in addition: an AMS \texttt{amsart} preamble that restates the article's own declarations and macros used by the retained lines, namely the environments \texttt{theorem}, \texttt{corollary} on separate counters, together with the standard packages those lines need.  The retained environments print in this document as Theorem 1, Corollary 1 in order of appearance, whereas the article prints the same items with the numbering of its own full document.  The complete native source of the article as distributed in the arXiv e-print for this exact version is bundled unchanged as \texttt{sources/208574/source0.tex}.
\begin{theorem}\label{thm:iso}
$M(\lambda)\cong M(\mu)$ as $V$-bimodules if and only if $\lambda-\mu\in\Fp$. In particular,
if $\lambda\in\Fp$ then
\[
\psi:V\longrightarrow M(\lambda),\qquad \psi(e_n)=m_{n-\lambda},
\]
is an isomorphism of bimodules, so $M(\lambda)\cong M(0)=V$.
\end{theorem}

\begin{corollary}\label{cor:notinFp}
If $\lambda\notin\Fp$ then $\Ht(V,M(\lambda))=0$.
\end{corollary}

\begin{align}
\varphi_{\delta g}(i,j)&=j\gamma_i+(j+d)\gamma_j-j\gamma_{i+j-1},\label{eq:cobV}\\
A(i,j,k)&=j\varphi(i+j-1,k)+k\varphi(i,j)-k\varphi(i,j+k-1)-(j+k-1+d)\varphi(j,k).
\label{eq:C1V}
\end{align}

\begin{theorem}\label{thm:p2}
Let $\operatorname{char}\kk=2$ and $V=\kk[x]/(x^2)$, so that $e_0=1$, $e_1=x$ and
\[
e_0\circ e_0=0,\qquad e_0\circ e_1=e_0,\qquad e_1\circ e_0=0,\qquad e_1\circ e_1=e_1 .
\]
If $\lambda\notin\mathbb{F}_2$ then $\Ht(V,M(\lambda))=0$. If $\lambda\in\mathbb{F}_2$ then
$M(\lambda)\cong V$ and
\[
\dim_\kk\Ht(V,M(\lambda))=\dim_\kk\Ht(V,V)=4 .
\]
More precisely, $\Ht_d(V,V)$ is $1$-dimensional for each $d\in\{-1,0,1,2\}$ and zero
otherwise, with representatives
\[
\begin{array}{lll}
d=2: &\Theta_1(e_i,e_j)=\delta_{i0}\delta_{j0}e_1, &\ \Theta_1(a,b)=a(0)\,b(0)\,x,\\[2pt]
d=1: &\Theta_2(e_i,e_j)=\delta_{i0}\delta_{j1}e_1, &\ \Theta_2(a,b)=a(0)\,b'(0)\,x,\\[2pt]
d=0: &\Theta_3(e_i,e_j)=\delta_{i1}\delta_{j0}e_0, &\ \Theta_3(a,b)=a'(0)\,b(0),\\[2pt]
d=-1:&\Theta_4(e_i,e_j)=\delta_{i1}\delta_{j1}e_0, &\ \Theta_4(a,b)=a'(0)\,b'(0).
\end{array}
\]
\end{theorem}

\begin{proof}
The first claim is Corollary~\ref{cor:notinFp}, whose proof is characteristic-free, and
$M(\lambda)\cong V$ is Theorem~\ref{thm:iso}. So it suffices to compute $\Ht(V,V)$, which we
do degree by degree in the $\mathbb{Z}$-graded truncated presentation. Here
$-2p+3=-1$ and $p=2$, so only $d\in\{-1,0,1,2\}$ admit nonzero cochains. Recall liveness
$0\le i+j-1+d\le1$, the window $0\le i+j+k-2+d\le1$, and that $\gamma_i$ is free iff
$0\le i+d\le1$; by \eqref{eq:cobV},
$\varphi_{\delta g}(i,j)=j\gamma_i+(j+d)\gamma_j-j\gamma_{i+j-1}$, and the last coefficient
in \eqref{eq:C1V} is $(j+k+d-1)$.

\textbf{$d=2$.} Liveness forces $i+j=0$: the only live pair is $(0,0)$, target $e_1$. No
$\gamma_i$ is free, so $\Bt_2=0$. The window forces $i+j+k\le1$. (C2) at $(0,0,1)$ and
$(0,1,0)$ are identities, and the triples $(0,0,0),(1,0,0)$ carry a factor $j=k=0$. For
(C1), $A(0,1,0)=\varphi(0,0)-(1+0+1)\varphi(1,0)=\varphi(0,0)$ (as $(1,0)$ is not live) and
$A(1,0,0)=-(0+0+1)\varphi(0,0)=\varphi(0,0)$ in characteristic $2$. So $\dim\Ht_2=1$.

\textbf{$d=1$.} Live pairs: $i+j\le1$, i.e.\ $(0,0),(0,1),(1,0)$ with targets $e_0,e_1,e_1$.
Only $\gamma_0$ is free and
\[
\varphi_{\delta g}(0,0)=\gamma_0,\qquad
\varphi_{\delta g}(0,1)=\gamma_0+(1+1)\gamma_1-\gamma_0=0,\qquad
\varphi_{\delta g}(1,0)=\gamma_0,
\]
using $\gamma_1=0$ and $2=0$; so $\Bt_1$ is spanned by $\varphi(0,0)=\varphi(1,0)=1$. The
window forces $i+j+k\le2$. From (C1) at $(0,1,0)$ and $(1,0,0)$,
\[
A(0,1,0)=\varphi(0,0)-(1+0)\varphi(1,0),\qquad A(1,0,0)=0,
\]
so $\varphi(0,0)=\varphi(1,0)$. At $(1,0,1)$ and $(0,1,1)$ one gets $-\varphi(0,1)$ and
$+\varphi(0,1)$, equal since $2=0$; all remaining triples and all of (C2) are identities.
Hence $\dim\Zt_1=2$, $\dim\Ht_1=1$, represented by $\varphi(0,1)=1$, i.e.\ by $\Theta_2$.

\textbf{$d=0$.} Live pairs: $1\le i+j\le2$, i.e.\ $(0,1),(1,0),(1,1)$ with targets
$e_0,e_0,e_1$. Both $\gamma_0,\gamma_1$ are free and
$\varphi_{\delta g}(i,j)=j(\gamma_i+\gamma_j+\gamma_{i+j-1})$, giving
$\varphi_{\delta g}(0,1)=\gamma_1$, $\varphi_{\delta g}(1,0)=0$,
$\varphi_{\delta g}(1,1)=\gamma_1$; so $\Bt_0$ is spanned by
$\varphi(0,1)=\varphi(1,1)=1$. The window forces $2\le i+j+k\le3$, and all instances of (C2)
are identities. For (C1),
\[
A(1,0,1)=\varphi(1,0)-\varphi(1,0)-(0+1-1)\varphi(0,1)=0,\quad
A(0,1,1)=\varphi(0,1)-(1+1-1)\varphi(1,1),
\]
so $\varphi(0,1)=\varphi(1,1)$; the triples $(1,1,0),(1,1,1)$ are symmetric in
$i\leftrightarrow j$. Hence $\dim\Zt_0=2$, $\dim\Ht_0=1$, represented by
$\varphi(1,0)=1$, i.e.\ by $\Theta_3$.

\textbf{$d=-1$.} Liveness forces $i+j=2$: the only live pair is $(1,1)$, target $e_0$. Only
$\gamma_1$ is free and
$\varphi_{\delta g}(1,1)=\gamma_1+(1-1)\gamma_1-\gamma_1=0$, so $\Bt_{-1}=0$. The window
forces $i+j+k=3$, leaving only $(1,1,1)$, where the coefficient $(j+k+d-1)$ vanishes and
$A(1,1,1)=\varphi(1,1)$ is symmetric in $i\leftrightarrow j$; (C2) is likewise an identity.
Hence $\dim\Ht_{-1}=1$, represented by $\Theta_4$.

Summing, $\dim\Ht(V,V)=4$.
\end{proof}

\end{document}
